\section{Praktikum}
\label{sec:praktikum-bab16}

\begin{praktikumbox}
\textbf{Judul:} Menganalisis, mengoptimasi, dan mengukur program pada
level instruksi.\\[0.3em]
\textbf{Sub-CPMK:} 4.3 \quad\textbullet\quad
\textbf{Alokasi waktu:} 150 menit \quad\textbullet\quad
\textbf{Bentuk:} kelompok (2--3 mahasiswa) dan laporan individu
\end{praktikumbox}

\subsection{Tujuan Praktikum}

Setelah menyelesaikan praktikum ini, mahasiswa mampu membongkar
program menjadi kode mesin, menemukan inefisiensi berdasarkan metrik,
menulis versi perbaikan, serta membuktikan perbaikannya melalui
pengukuran yang adil.

\subsection{Alat dan Bahan}

\begin{itemize}
    \item Komputer dengan DOSBox dan MS-DEBUG, atau Emu8086.
    \item TASM dan TLINK (atau GUI Turbo Assembler).
    \item Berkas \file{naif.asm}, \file{teroptimasi.asm}, dan
          \file{teroptimasi-unroll.asm} dari \file{code/asm/}.
    \item Lembar kerja pengukuran dan formulir jurnal praktikum.
\end{itemize}

\subsection{Langkah Kerja}

\begin{enumerate}
    \item \textbf{Merakit dan menjalankan tiga versi program.}
    \begin{enumerate}
        \item Rakit ketiga berkas dengan perintah \file{tasm} dan
              \file{tlink} di dalam DOSBox.
        \item Jalankan tiap versi, lalu periksa isi variabel
              \code{hasil} melalui jendela memori Emu8086 atau perintah
              \code{D} pada DEBUG untuk memastikan ketiganya
              menghasilkan \hex{1130}.
        \item Catat panjang berkas \file{.exe} masing-masing versi.
    \end{enumerate}

    \item \textbf{Membongkar kode mesin.}
    \begin{enumerate}
        \item Muat program dengan perintah \code{N} dan \code{L} pada
              DEBUG, lalu bongkar mulai titik masuk program dengan
              perintah \code{U}.
        \item Catat alamat, byte mesin, dan \eng{mnemonic} setiap
              instruksi pada lembar kerja.
        \item Tandai setiap pola yang telah dipelajari pada
              Tabel~\ref{tab:pola-disassembly}: perulangan,
              inisialisasi segmen, dan sebagainya.
        \item Ulangi untuk versi teroptimasi, lalu susun kedua hasil
              secara berdampingan seperti pada
              Gambar~\ref{fig:disasm-bandingan}.
    \end{enumerate}

    \item \textbf{Mengidentifikasi inefisiensi.}
    \begin{enumerate}
        \item Tandai instruksi pada versi naif yang memakai akses
              memori berulang, memuat konstanta panjang, atau
              memakai \asm{MUL}/\asm{DIV}.
        \item Hitung perkiraan siklus bagian yang diukur dengan
              bantuan Tabel~\ref{tab:siklus-8086-perkiraan}.
        \item Tuliskan usulan perbaikan beserta alasan dan syarat
              berlakunya.
    \end{enumerate}

    \item \textbf{Menulis versi perbaikan sendiri.}
    \begin{enumerate}
        \item Salin \file{naif.asm} menjadi \file{optimasi-anda.asm}.
        \item Terapkan minimal tiga teknik dari
              Tabel~\ref{tab:strength-reduction}.
        \item Pastikan keluaran tetap sama: variabel \code{hasil}
              harus tetap bernilai \hex{1130}.
    \end{enumerate}

    \item \textbf{Mengukur secara adil.}
    \begin{enumerate}
        \item Tambahkan penghitung iterasi dan pencatat waktu
              \code{INT 1Ah} fungsi \hex{00} seperti pada
              Subbab~\ref{sec:optimasi-instruksi}.
        \item Ukur ketiga versi (naif, teroptimasi, dan
              \file{optimasi-anda.asm}) dengan larik dan jumlah
              iterasi yang sama.
        \item Ulangi setiap pengukuran lima kali, buang percobaan
              pertama, lalu hitung nilai tengahnya.
        \item Isi Tabel~\ref{tab:hasil-pengukuran-16} dengan hasil
              pengukuran Anda, bukan dengan angka perkiraan.
    \end{enumerate}
\end{enumerate}

\subsection{Hasil yang Dikumpulkan}

\begin{enumerate}
    \item Berkas \file{optimasi-anda.asm} yang dapat dirakit tanpa
          kesalahan.
    \item Tabel perbandingan hasil \eng{disassembly} versi naif dan
          versi perbaikan.
    \item Tabel hasil pengukuran sebelum dan sesudah, lengkap dengan
          perkiraan siklus dan jumlah akses memori.
    \item Analisis singkat: teknik mana yang memberi perbaikan
          terbesar, dan teknik mana yang menambah ukuran kode.
\end{enumerate}

\subsection{Studi Kasus dan Presentasi Proyek Akhir}

Praktikum terakhir semester ini sekaligus menjadi tempat studi kasus
dan presentasi proyek akhir. Setiap kelompok memilih satu program
dari bab-bab sebelumnya, lalu menjalankan seluruh alur kerja bab ini
terhadap program tersebut.

\begin{studikasusbox}
\textbf{Studi kasus: ``Program mahasiswa yang lambat''.} Sebuah
program pengolah larik buatan mahasiswa memerlukan waktu yang jauh
lebih lama daripada perkiraan pemiliknya. Program itu benar dan
keluarannya tepat, sehingga tidak ada kesalahan yang dapat ditemukan
dengan \eng{debugging} biasa. Tugas kelompok adalah menemukan
penyebabnya dengan perintah \code{U}, menunjukkannya melalui metrik,
memperbaikinya, dan membuktikan perbaikan itu dengan pengukuran.

\vspace{0.4em}
\textbf{Panduan singkat presentasi (7 menit + 3 menit tanya jawab):}
\begin{enumerate}
    \item Tampilkan potongan hasil \eng{disassembly} sebelum dan
          sesudah, dan tunjuk satu instruksi yang menjadi biang
          kelambatan (2 menit).
    \item Jelaskan teknik optimasi yang dipilih beserta syarat
          berlakunya (2 menit).
    \item Tampilkan tabel pengukuran sebelum dan sesudah; sebutkan
          metrik yang dipakai dan asumsi pengukurannya (2 menit).
    \item Nyatakan batas pekerjaan: apa yang \emph{tidak} berubah
          dan risiko apa yang masih tersisa (1 menit).
\end{enumerate}
Semua anggota kelompok harus mengambil bagian berbicara, dan setiap
klaim perbaikan wajib disertai angka.
\end{studikasusbox}

\begin{asesmenbox}
\textbf{Asesmen praktikum 15 (bobot 15\% dari komponen praktikum).}

\begin{tabular}{@{}p{6.0cm}cp{4.8cm}@{}}
    \toprule
    \textbf{Indikator} & \textbf{Bobot} & \textbf{Kriteria} \\
    \midrule
    Ketepatan hasil \eng{disassembly} & 25\% &
    Alamat, byte mesin, dan \eng{mnemonic} tercatat benar serta pola
    dikenali dengan tepat \\
    \addlinespace
    Kualitas identifikasi inefisiensi & 20\% &
    Minimal tiga inefisiensi ditemukan dan dikaitkan dengan metrik \\
    \addlinespace
    Kebenaran versi perbaikan & 25\% &
    Program berjalan, hasil \hex{1130} tetap, dan minimal tiga
    teknik optimasi diterapkan sesuai syaratnya \\
    \addlinespace
    Keabsahan pengukuran & 20\% &
    Masukan dan jumlah iterasi sama; pengukuran diulang dan nilai
    tengah dilaporkan \\
    \addlinespace
    Kerapian laporan dan jurnal & 10\% &
    Langkah kerja dapat diulang oleh pembaca lain \\
    \bottomrule
\end{tabular}
\end{asesmenbox}

\begin{asesmenbox}
\textbf{Rubrik presentasi proyek akhir (bobot 10\% dari komponen
praktikum).}

\begin{tabular}{@{}p{4.2cm}p{1.6cm}p{5.0cm}@{}}
    \toprule
    \textbf{Aspek} & \textbf{Bobot} & \textbf{Kriteria penilaian} \\
    \midrule
    Penguasaan hasil \eng{disassembly} & 25\% &
    Dapat menjelaskan makna byte dan pola tiap instruksi yang
    ditampilkan \\
    \addlinespace
    Argumentasi optimasi & 25\% &
    Setiap perubahan disertai alasan, syarat berlaku, dan
    kesadaran akan trade-off \\
    \addlinespace
    Bukti pengukuran & 25\% &
    Tabel sebelum/sesudah lengkap; keadilan pengukuran dijelaskan \\
    \addlinespace
    Penyajian dan tanya jawab & 25\% &
    Sistematis, waktu terpenuhi, seluruh anggota berbicara, jawaban
    tepat sasaran \\
    \bottomrule
\end{tabular}
\end{asesmenbox}
